\chapter{Photonic Emission Channels and Phase Crosstalk}
\label{app:photonic_emission}
\addtocontents{toc}{\protect\setcounter{tocdepth}{0}}

\noindent\textbf{Abstract:} In Tension Field Theory (TFT), a photon is defined as a coherent release of substrate tension along a discrete electromagnetic channel. This appendix formalizes the mechanics of photon emission from atomic orbital drops, proving the existence of picometer-scale emission bottlenecks. Furthermore, it defines the phase exchange budget for independent photons crossing paths via temporary strain-field overlap on the SC Substrate.

\section{The Mechanics of Photon Emission}
Standard physics models photon emission as a probabilistic point-particle creation. TFT models it as a mechanical extrusion of structural tension. Photon emission is the release of structural tension through a highly localized micro-bottleneck in the FCC Soft Floor.

The energy of the emitted photon is the volume integral of the released tension over the geometric bounds of the emission channel:
\begin{equation}
    E_{\gamma} = \int \delta\tau \, \mathrm{d}V \approx \tau_{\mathrm{bulk}} \cdot A_{\mathrm{chan}} \cdot \Delta r
\end{equation}

\noindent\textbf{Where:}
\begin{itemize}
    \item $\tau_{\mathrm{bulk}}$: The three-dimensional tension density of the local field (N/m$^2$). The fundamental macroscopic tension is defined by the lattice density $\tau_{\mathrm{bulk}} = \rho c_{\mathrm{E}}^2$. Substituting the geometric definition of the quantum action lever ($\hbar = \rho c_{\mathrm{E}} r^4$), the operative tension density at a given effective radius $r_{\mathrm{eff}}$ scales strictly as:
    \begin{equation}
        \tau_{\mathrm{bulk}} = \frac{\hbar c_{\mathrm{E}}}{r_{\mathrm{eff}}^4}
    \end{equation}
    \item $A_{\mathrm{chan}}$: The cross-sectional area of the micro-emission channel.
    \item $\Delta r$: The physical drop length through the lattice, defined by the absolute spatial difference between the starting shell $r_n$ and ending shell $r_m$: $\Delta r = |r_n - r_m|$.
    \item $r_{\mathrm{eff}}$: The operative geometric mean radius of the atomic drop: $r_{\mathrm{eff}} = \sqrt{r_n r_m}$.
\end{itemize}

\subsection{Worked Hydrogen Examples}
By tracing the exact geometric path of the electron drop, we can calculate the true cross-section of the emission channel ($A_{\mathrm{chan}} = \frac{|\Delta E|}{\tau_{\mathrm{bulk}} \Delta r}$).

\vspace{0.2cm}
\noindent \textbf{1. Lyman--$\alpha$ ($2 \to 1$):}
\begin{itemize}
    \item \textbf{$\Delta E$:} $\approx 10.204$ eV \quad ($1.6349 \times 10^{-18}$ J)
    \item \textbf{Geometry:} $r_2 = 2.1167 \times 10^{-10}$ m, \quad $r_1 = 5.2918 \times 10^{-11}$ m
    \item \textbf{Drop Parameters:} $\Delta r = 1.5875 \times 10^{-10}$ m, \quad $r_{\mathrm{eff}} = 1.0584 \times 10^{-10}$ m
    \item \textbf{Local Tension:} $\tau_{\mathrm{bulk}} \approx 2.5220 \times 10^{14}$ N/m$^2$
    \item \textbf{Channel Area:} $A_{\mathrm{chan}} \approx 4.0835 \times 10^{-23}$ m$^2$
    \item \textbf{Resulting Radius:} $r_{\mathrm{chan}} \approx 3.6057$ pm \quad ($3.6057 \times 10^{-12}$ m)
\end{itemize}

\vspace{0.2cm}
\noindent \textbf{2. Balmer--$\alpha$ ($3 \to 2$):}
\begin{itemize}
    \item \textbf{$\Delta E$:} $\approx 1.8897$ eV \quad ($3.0276 \times 10^{-19}$ J)
    \item \textbf{Geometry:} $r_3 = 4.7626 \times 10^{-10}$ m, \quad $r_2 = 2.1167 \times 10^{-10}$ m
    \item \textbf{Drop Parameters:} $\Delta r = 2.6459 \times 10^{-10}$ m, \quad $r_{\mathrm{eff}} = 3.1751 \times 10^{-10}$ m
    \item \textbf{Local Tension:} $\tau_{\mathrm{bulk}} \approx 3.1112 \times 10^{12}$ N/m$^2$
    \item \textbf{Channel Area:} $A_{\mathrm{chan}} \approx 3.6779 \times 10^{-22}$ m$^2$
    \item \textbf{Resulting Radius:} $r_{\mathrm{chan}} \approx 10.820$ pm \quad ($1.0820 \times 10^{-11}$ m)
\end{itemize}

\noindent\textbf{Conclusion:} Both atomic transitions prove mathematically that photons are extruded through localized, picometer-scale channels ($\approx 10^{-12}$ m), which are vastly smaller than the macroscopic atomic orbital radii ($\approx 10^{-10}$ m). The released energy does not occupy the full orbital shell; it is squeezed out through a highly pressurized micro-fracture.

\section{The Phase Crosstalk Budget (Shared Strain Bottleneck)}
When two independent photons cross paths in free space at distances vastly greater than their channel radii, they exchange nothing.

However, if their trajectories are engineered to overlap tightly (such as in an evanescent fiber coupler), their vertical Tension Columns pass close enough that their respective strain fields on the SC Hard Floor temporarily overlap. This creates a Shared Strain Bottleneck, $A_{\mathrm{overlap}}$.

Because they temporarily share a localized patch of the phase-locked SC substrate, the torsion of Column A mechanically interacts with the torsion of Column B. They do not exchange transverse kinetic energy (which requires field fracture); they strictly exchange Phase Holonomy, $\Delta\varphi$.

This mechanical phase exchange is governed by the total volumetric action of the overlapping strain fields, bounded by:
\begin{equation}
    \Delta\varphi \leq \eta_{\varphi} \cdot \left(\frac{c_{\mathrm{E}}}{r^3}\right) \cdot A_{\mathrm{overlap}} \cdot T
\end{equation}

\noindent\textbf{Where:}
\begin{itemize}
    \item $r$ is the mode radius.
    \item $A_{\mathrm{overlap}}$ is the transient bottleneck area of the crossing Hard Floor strain fields.
    \item $T$ is the operational dwell time.
    \item $\eta_{\varphi}$ is the dimensionless Geometric Coupling Efficiency ($0 < \eta_{\varphi} \ll 1$). It is an empirically calibrated constant that accounts for spatial mismatch, scattering losses, and the imperfect alignment of the transient strain fields during the crossing event.
\end{itemize}

\section{Phase-Crosstalk Predictions and Falsification}
\label{sec:AG_predictions_falsification}

When two coherent photons cross within a sufficiently confined
region, TFT proposes that their overlapping Active-Vacuum strain
fields may exchange phase information without requiring a
corresponding transfer of carrier energy. This fleeting
mechanical interaction is described by the Shared Strain
Bottleneck, whose geometry and interaction duration constrain
the possible phase exchange.

\subsection*{TFT Prediction: The Phase-Crosstalk Budget}

Within the proposed phase-transfer model, the additional
phase magnitude is bounded by
\[
0\leq|\Delta\varphi|
\leq
\eta_{\varphi}
\frac{c_{\mathrm E}}{r^3}
A_{\mathrm{overlap}}T.
\]
Here $\eta_{\varphi}$ represents the geometric coupling
efficiency, $r$ the mode radius, $A_{\mathrm{overlap}}$
the shared strain area, and $T$ the interaction dwell time.

When the remaining variables are held fixed, the upper
bound scales linearly with overlap area and dwell time,
and inversely with the cube of the mode radius.
These relations define an upper bound, not a compulsory
phase shift. A smaller mode radius may also alter the
overlap and coupling efficiency, whose physical responses
must be considered together.

\subsection*{Falsification Criteria}

\noindent\textbf{The Phase-Budget Limit:}
If a measured interaction-induced phase shift exceeds
the maximum permitted budget, with the relevant geometry,
interaction time, and coupling efficiency independently
constrained, the proposed phase-crosstalk closure is
falsified. An unrestricted coupling coefficient cannot
be readjusted after each measurement to preserve the bound.

Conversely, a null phase shift or an inverse-square
distance dependence does not invalidate an upper-bound
relation, provided the measured response remains within
its allowed range. The present construction does not
yet determine a universal numerical phase cap.

\medskip
\noindent\textbf{Phase Exchange versus Energy Transfer:}
The proposed phase-only channel must also be distinguished
from genuinely inelastic photon interactions. A measured
frequency change does not automatically contradict TFT,
since material coupling or other physical channels may
transfer energy. However, if a completed microscopic
phase-only interaction predicts negligible energy transfer
under specified conditions, reproducible energy exchange
attributable to that same interaction would falsify
the phase-only realization.

\medskip
The significance of this test lies in separating the
geometry that permits an interaction from the mechanics
that determines its strength. The bottleneck constrains
the available phase budget; the actual exchange must
follow from a physical coupling law.

\medskip
\noindent\textbf{The overlap constrains the phase budget;
the interaction determines the actual exchange, and
observation decides whether the mechanism survives.}